% 1_motivation.tex
\section{Motivation}
\label{sec:ch1_motivation}

% Paragraph 1: The orchestration challenge (unified)
\ac{6G} networks demand orchestration systems that manage hundreds of distributed nodes~\cite{6g_roaming, intent_driven_ran2021}. These systems must support services requiring sub-millisecond latency and strict \ac{QoS} guarantees. Cloud-native architectures require coordination of thousands of microservices distributed across far-edge devices, near-edge nodes, and data centers~\cite{6g_struggles_orchestration2025}. Traditional centralized orchestration platforms struggle with the latency and scalability requirements of this distributed coordination. As wireless systems evolve from monolithic to microservice-based architectures, they introduce dependencies, unpredictable demand fluctuations, and frequent container migrations~\cite{10029931, ouyang2023dynamic}. This management complexity causes severe resource imbalances. Analysis of production cloud workloads reveals that average machine utilization remains below 40\%, and most service instances consume less than 10\% of their reserved \ac{CPU} resources~\cite{8258257}. These inefficiencies span spatial and temporal dimensions that current orchestration methods struggle to address.

% Paragraph 2: Traditional approaches insufficient (unified)
Three fundamental limitations of traditional orchestration approaches contribute to these inefficiencies. First, they operate reactively rather than predictively~\cite{aloqaily2023metaverse}. This creates a dilemma where allocating excess capacity wastes resources while allocating insufficient capacity degrades service quality. Second, manual policies struggle to adapt to diverse workload patterns~\cite{chen2018auto}. Services follow daily periodicities driven by user activity, while batch processing generates irregular bursts determined by job submissions~\cite{calzarossa2016workload}. Without adaptive intelligence, static configurations perform poorly across such heterogeneous workloads. Third, centralized decision-making architectures introduce communication latency that affects responsiveness in edge deployment scenarios. They also create single points of failure that cascade across distributed infrastructures~\cite{6g_struggles_orchestration2025}. These limitations are interconnected constraints rather than independent deficiencies. Together, they limit the ability of traditional methods to achieve the efficiency, scalability, and reliability demanded by \ac{6G} networks.

% Paragraph 3: What intelligent orchestration requires (unified vision)
Addressing these limitations requires intelligent orchestration systems that learn from operational data, adapt to diverse application behaviors, and coordinate decisions across distributed infrastructure without human intervention. To overcome the reactive nature, adaptation challenges, and coordination complexity identified above, such systems must fulfill three capabilities: (i) predict resource demands before they materialize, enabling proactive allocation that reduces reactive waste, (ii) generalize learned behaviors across heterogeneous application types, avoiding unsustainable per-service customization overhead, and (iii) reason autonomously over complex multi-dimensional constraints to coordinate orchestration actions across distributed nodes. Traditional approaches separate prediction from execution, require manual tuning for each application, and operate through reactive policies. In contrast, intelligent orchestration integrates learning, adaptation, and autonomous coordination into unified frameworks.

% Paragraph 4: Three interdependent dimensions (natural introduction - merged with ML justification)
These capabilities align with the strengths of machine learning and artificial intelligence: pattern recognition from data, generalization across contexts, and multi-dimensional reasoning. However, applying these techniques to \ac{6G} orchestration introduces specific technical challenges. This dissertation structures these challenges as three interdependent dimensions of the resource management lifecycle. \textit{Resource efficiency} benefits from predictions occurring close to where orchestration decisions happen, at resource-constrained edge devices where centralized approaches can introduce significant communication latency. \textit{Operational scalability} requires that solutions serve diverse application types without per-service customization. Maintaining hundreds of specialized models creates substantial operational overhead. \textit{Autonomous management} requires systems that reason over complex constraints and coordinate actions across distributed nodes without human intervention. Manual decision-making cannot scale to the complexity of managing thousands of interdependent services. Edge deployment enables timely predictions, generalization enables scalability across applications, and autonomous reasoning closes the management loop by executing coordinated orchestration actions. The following sections examine each dimension in depth.

% Paragraph 5: Resource efficiency dimension - deep dive
The resource efficiency dimension addresses a practical constraint: deploying prediction models at edge devices reduces network delays but requires operating within limited computational resources. Existing learning-based forecasting approaches are computationally intensive, requiring memory and processing capacity that exceed typical edge device constraints~\cite{kong2022edge, 8334540, 9500858}. State-of-the-art methods achieve high accuracy through large model architectures suitable for centralized cloud deployment but infeasible for distributed edge execution~\cite{meuser2024revisiting}. This creates a critical tension where accurate workload prediction requires sophisticated pattern learning while edge constraints demand lightweight computation. Centralized prediction can diminish the benefits of proactive resource allocation because communication latency between cloud forecasting services and edge orchestration components introduces delays that diminish the benefits of proactive allocation. Therefore, achieving resource efficiency requires rethinking model design to maintain prediction accuracy within edge deployment constraints rather than simply compressing existing architectures.

% Paragraph 6: Operational scalability dimension - deep dive
Solving the hardware constraint with lightweight models introduces a new challenge: operational scalability. While lightweight architectures can fit on edge devices, they typically rely on learning specific local patterns, such as the periodicity of a single service. As a result, they struggle to generalize across the heterogeneous behaviors of thousands of different microservices~\cite{zou2024generalization}. Deploying these models at scale would require training and maintaining a unique instance for every node and service in the network. This per-application customization diminishes the benefits of lightweight prediction because the operational cost of managing thousands of specialized models can exceed the cost of manual orchestration. Achieving operational scalability thus requires discovering fundamental workload patterns that transcend application-specific implementations, enabling single models to serve multiple service types by learning which behavioral patterns exist and which apply to each application context.

% Paragraph 7: Autonomous management dimension - deep dive (expanded)
Precise and scalable predictions have limited value without a system to act on them. The autonomous management dimension addresses the fundamental limitation that orchestration complexity exceeds human decision-making capacity. Managing hundreds of heterogeneous services requires continuous reasoning over multi-dimensional constraints: infrastructure topology, workload dynamics, \ac{QoS} requirements, resource availability, and service dependencies. Each constraint affects orchestration decisions differently, and all evolve dynamically as conditions change. This multi-dimensional optimization problem involves thousands of interdependent parameters. These parameters evolve dynamically as infrastructure states change, workloads fluctuate, and failures occur. Human operators face significant challenges processing this complexity at the speed and scale demanded by \ac{6G} networks. Consequently, \ac{6G} orchestration requires autonomous systems capable of interpreting high-level management intentions, reasoning over complex infrastructure states, planning multi-step orchestration strategies, executing coordinated actions across distributed nodes, and adapting to unexpected failures or demand shifts without human intervention. The fundamental challenge is achieving true zero-touch orchestration where systems operate autonomously without requiring manual configuration, human oversight, or intervention during normal operations. Such systems must also scale hierarchically to manage geographically distributed infrastructure.

% Paragraph 8: Dissertation approach (three integrated frameworks)
Building on this analysis, this dissertation addresses these three interdependent dimensions through complementary frameworks that combine learning-based intelligence with orchestration-specific design principles. For resource efficiency, lightweight forecasting models are designed from the ground up to use domain-specific patterns inherent in microservice workloads, enabling accurate prediction within edge device constraints. For operational scalability, pattern discovery methods identify fundamental workload behaviors that transcend application-specific details, enabling single models to serve diverse applications without per-service retraining. For autonomous management, multi-agent frameworks decompose orchestration complexity into specialized reasoning roles that coordinate through structured workflows, enabling zero-touch operation where systems make and execute decisions autonomously.

% Paragraph 9: Integration and interdependence
The three dissertation contributions form a complementary hierarchy of intelligence where each layer addresses a scope of the orchestration problem that the others cannot. \textbf{AERO (Chapter~\ref{ch:adaptive-orchestration})} serves as the \textit{local specialist}, running directly on edge devices to manage immediate resource acceptance and job scheduling with minimal overhead. \textbf{OmniFORE (Chapter~\ref{ch:attention-driven-prediction})} provides \textit{network-level generalization}, predicting for many applications using one model instead of one per service. \textbf{AgentEdge (Chapter~\ref{ch:distributed-agentic-ai})} functions as the \textit{autonomous executive}, turning predictions into coordinated orchestration decisions across distributed infrastructure without requiring human operators. Together, these frameworks target efficiency at the device level, scalability across network regions, and autonomous operation across distributed infrastructure.
